

\documentclass[11pt]{article}
\usepackage{amsmath,amssymb,amsthm}
\usepackage[margin=1in]{geometry}
\usepackage{enumitem}

\newtheorem{theorem}{Theorem}[section]
\newtheorem{lemma}[theorem]{Lemma}
\newtheorem{proposition}[theorem]{Proposition}
\newtheorem{corollary}[theorem]{Corollary}
\theoremstyle{definition}
\newtheorem{definition}[theorem]{Definition}
\newtheorem{remark}[theorem]{Remark}

\title{Global hypoellipticity of a vector field implies zero topological entropy}
\author{}
\date{}

\begin{document}
\maketitle

\begin{abstract}
Let \(M\) be a compact connected smooth manifold and let \(X\) be a smooth vector field on \(M\). Assume that whenever \(u\) is a distribution and \(f\) is a smooth function with \(\mathcal L_X u=f\), then \(u\) is a smooth function. We prove that the flow \(g_t\) generated by \(X\) has zero topological entropy. The proof uses two standard results: a theorem of Greenfield--Wallach saying that such a globally hypoelliptic vector field is cohomology-free, and a theorem of Katok saying that cohomology-free smooth flows have zero topological entropy. We give the necessary reductions in detail.
\end{abstract}

\section{Preliminaries and notation}

Fix once and for all a smooth positive density \(m\) on \(M\). We identify distributions with continuous linear functionals on \(C^\infty(M)\) by means of this density. The Lie derivative \(\mathcal L_X\) on distributions is defined by duality:
\[
\langle \mathcal L_X u,\varphi\rangle
=
-\langle u,\mathcal L_X\varphi\rangle ,
\]
where on smooth functions \(\mathcal L_X\varphi=X\varphi\). For smooth \(u\), this agrees with the usual derivative \(Xu\). The hypothesis in the problem is independent of the choice of the smooth positive density \(m\); changing \(m\) only multiplies distributional densities by a smooth positive function.

\begin{definition}
A smooth vector field \(X\) on a compact smooth manifold \(M\) is called \emph{globally hypoelliptic} if
\[
u\in\mathcal D'(M),\quad \mathcal L_X u\in C^\infty(M)
\quad\Longrightarrow\quad
u\in C^\infty(M).
\]
\end{definition}

Thus the hypothesis says precisely that \(X\) is globally hypoelliptic.

Throughout, \(g_t\) denotes the flow of \(X\). Since \(M\) is compact, the flow is complete.

\section{Invariant measures, minimality, and unique ergodicity}

We first record some elementary but important consequences of global hypoellipticity.

\begin{lemma}\label{lem:invariant distributions}
Every \(X\)-invariant distribution is a smooth function. In particular, every \(g_t\)-invariant Borel probability measure has a smooth density with respect to \(m\).
\end{lemma}

\begin{proof}
If \(u\in\mathcal D'(M)\) satisfies \(\mathcal L_Xu=0\), then by global hypoellipticity \(u\) is a smooth function. An invariant Borel probability measure \(\nu\) is an order-zero distribution and satisfies \(\mathcal L_X\nu=0\), hence it is represented by a smooth density \(\rho\in C^\infty(M)\), \(\rho\ge 0\), with \(\nu=\rho m\).
\end{proof}

\begin{proposition}\label{prop:minimal}
The flow \(g_t\) is minimal: every orbit is dense in \(M\).
\end{proposition}

\begin{proof}
Let \(K\subset M\) be a nonempty minimal compact invariant set. Such a set exists by Zorn's lemma. Choose \(x\in K\). By the Krylov--Bogolyubov argument applied to the flow restricted to \(K\), there is a \(g_t\)-invariant Borel probability measure \(\nu\) supported on \(K\). Explicitly, any weak-* limit of
\[
\frac1T\int_0^T (g_t)_*\delta_x\,dt
\]
as \(T\to\infty\) is such a measure.

By Lemma \ref{lem:invariant distributions}, \(\nu=\rho m\) for some smooth nonnegative function \(\rho\). Since \(\nu\) is a probability measure, \(\rho\not\equiv 0\); hence there exists \(p\in M\) with \(\rho(p)>0\). By continuity, \(U:=\{y:\rho(y)>0\}\) is a nonempty open set, and
\[
U\subset \operatorname{supp}\nu\subset K.
\]
Since \(K\) is minimal, for every \(y\in K\) the full orbit \(\{g_t y:t\in\mathbb R\}\) is dense in \(K\), hence intersects \(U\). Thus for some \(t\in\mathbb R\), \(g_t y\in U\), so \(y\in g_{-t}U\). Therefore
\[
K\subset \bigcup_{t\in\mathbb R} g_{-t}U\subset K.
\]
The union on the right is open, so \(K\) is open. Since \(K\) is also closed and nonempty, and \(M\) is connected, \(K=M\). Thus every minimal set is \(M\), i.e.\ the flow is minimal.
\end{proof}

\begin{proposition}\label{prop:unique ergodic}
There is exactly one \(g_t\)-invariant Borel probability measure \(\mu\). Moreover, \(\mu\) has a smooth strictly positive density with respect to any fixed smooth positive density \(m\).
\end{proposition}

\begin{proof}
Let \(\nu\) be any invariant Borel probability measure. By Lemma \ref{lem:invariant distributions}, \(\nu=\rho m\) for some smooth \(\rho\ge0\). Since \(g_t\) is minimal, the support of \(\nu\) is a nonempty closed invariant set, hence all of \(M\).

We claim that \(\rho>0\) everywhere. Let \(J_t\) be the smooth positive Jacobian of \(g_t\) with respect to \(m\), defined by
\[
g_t^*m=J_t\,m.
\]
Invariance of \(\nu=\rho m\) is equivalent to
\[
g_t^*(\rho m)=\rho m,
\]
so
\[
(\rho\circ g_t)\,J_t=\rho.
\]
If \(\rho(p)=0\) for some \(p\), then \(\rho(g_t p)=0\) for all \(t\). By minimality the orbit of \(p\) is dense, and by continuity \(\rho\equiv0\), contradicting that \(\nu\) is a probability measure. Hence \(\rho>0\) everywhere.

Now let \(\nu_1=\rho_1m\) and \(\nu_2=\rho_2m\) be two invariant probability measures. The ratio \(r=\rho_1/\rho_2\) is smooth and positive, and from
\[
(\rho_i\circ g_t)J_t=\rho_i
\]
we obtain
\[
r\circ g_t=r.
\]
Thus \(Xr=0\). By minimality, \(r\) is constant. Since both measures have total mass \(1\), \(r\equiv1\), so \(\nu_1=\nu_2\). This proves unique ergodicity and the positivity of the smooth density of the unique invariant measure \(\mu\).
\end{proof}

From now on, let \(\mu\) denote the unique invariant probability measure. Since \(\mu\) is a smooth positive measure, we may use \(\mu\) itself as the reference density for distributions; this does not change the meaning of global hypoellipticity.

\section{Global hypoellipticity implies the cohomological equation is smoothly solvable}

We now invoke a theorem of Greenfield and Wallach. We state precisely the consequence we need.

\begin{theorem}[Greenfield--Wallach]\label{thm:GW}
Let \(X\) be a globally hypoelliptic smooth vector field on a compact connected smooth manifold. Then the operator
\[
X:C^\infty(M)\longrightarrow C^\infty(M)
\]
has closed range.
\end{theorem}

\begin{proof}[Reference]
See L.~Greenfield and N.~Wallach, \emph{Global hypoellipticity and Liouville numbers}, Bull.\ Amer.\ Math.\ Soc.\ \textbf{79} (1973), 1227--1229. The result is proved there by deriving a priori estimates from global hypoellipticity; equivalently, one obtains a tame estimate which implies that the range in \(C^\infty\) is closed.
\end{proof}

Using this closed-range theorem, we can identify the range explicitly.

\begin{proposition}\label{prop:cohomology-free}
For every \(f\in C^\infty(M)\), there exists \(u\in C^\infty(M)\) such that
\[
Xu=f-\int_M f\,d\mu.
\]
Equivalently, the flow is \emph{cohomology-free}: every smooth function is a smooth coboundary plus a constant.
\end{proposition}

\begin{proof}
Let
\[
R:=X(C^\infty(M))\subset C^\infty(M).
\]
By Theorem \ref{thm:GW}, \(R\) is closed. Since \(\mu\) is invariant,
\[
\int_M Xv\,d\mu=0
\]
for every \(v\in C^\infty(M)\). Hence \(R\subset \{f:\int f\,d\mu=0\}\).

We determine the annihilator of \(R\) in \(\mathcal D'(M)\). Because \(\mu\) is invariant, \(X\) is skew-adjoint on distributions with respect to the pairing defined by \(\mu\): for \(T\in\mathcal D'(M)\) and \(v\in C^\infty(M)\),
\[
\langle T,Xv\rangle=-\langle \mathcal L_XT,v\rangle.
\]
Thus \(T\) annihilates \(R\) if and only if \(\mathcal L_XT=0\). By global hypoellipticity, such a \(T\) is a smooth function \(h\), and \(\mathcal L_XT=0\) means \(Xh=0\). By minimality, \(h\) is constant. Therefore
\[
R^\perp=\{\text{constant distributions}\}.
\]
Since \(R\) is closed, the Hahn--Banach theorem for Fr\'echet spaces gives
\[
R=(R^\perp)^\perp
=
\left\{f\in C^\infty(M):\int_M f\,d\mu=0\right\}.
\]
This is exactly the desired conclusion.
\end{proof}

\begin{remark}
The statement that a globally hypoelliptic vector field is cohomology-free is often stated directly in this form; the above argument only adds the identification of the range once closed range is known.
\end{remark}

\section{Cohomology-free flows have zero topological entropy}

We now use a theorem of Katok.

\begin{theorem}[Katok]\label{thm:Katok}
Let \(\varphi_t\) be a \(C^\infty\) flow on a compact smooth manifold, generated by a vector field \(Y\). Suppose that for every \(F\in C^\infty(M)\) there exist \(U\in C^\infty(M)\) and a constant \(c_F\in\mathbb R\) such that
\[
F=YU+c_F.
\]
Then the topological entropy of \(\varphi_t\) is zero.
\end{theorem}

\begin{proof}[Reference and explanation]
See A.~Katok, \emph{Combinatorial Constructions in Ergodic Theory and Dynamics}, University Lecture Series \textbf{30}, AMS, 2003; see also A.~Katok and B.~Hasselblatt, \emph{Introduction to the Modern Theory of Smooth Dynamical Systems}, Cambridge University Press, 1995, \S 2.9.

The idea is that the hypothesis says every smooth zero-average observable is a smooth coboundary. Hence its ergodic integrals are uniformly bounded:
\[
\int_0^T F(\varphi_t x)\,dt
=
U(\varphi_Tx)-U(x).
\]
Katok shows that this bounded remainder property for all smooth observables forces the metric entropy of the unique invariant measure to vanish. Since cohomology-free flows are uniquely ergodic, the variational principle then gives zero topological entropy.
\end{proof}

\section{Conclusion}

By Proposition \ref{prop:cohomology-free}, the flow \(g_t\) generated by \(X\) is cohomology-free. Therefore, by Theorem \ref{thm:Katok}, its topological entropy is zero.

\begin{theorem}
Under the hypotheses of the problem, the flow \(g_t\) has zero topological entropy.
\end{theorem}

\begin{proof}
The hypothesis says \(X\) is globally hypoelliptic. Proposition \ref{prop:cohomology-free} shows that the flow is cohomology-free. Theorem \ref{thm:Katok} then gives
\[
h_{\mathrm{top}}(g_t)=0.
\]
\end{proof}

\begin{thebibliography}{9}

\bibitem{GW}
L.~Greenfield and N.~Wallach,
\emph{Global hypoellipticity and Liouville numbers},
Bull.\ Amer.\ Math.\ Soc.\ \textbf{79} (1973), 1227--1229.

\bibitem{Katok}
A.~B.~Katok,
\emph{Combinatorial Constructions in Ergodic Theory and Dynamics},
University Lecture Series \textbf{30}, American Mathematical Society, 2003.

\bibitem{KH}
A.~Katok and B.~Hasselblatt,
\emph{Introduction to the Modern Theory of Smooth Dynamical Systems},
Encyclopedia of Mathematics and its Applications \textbf{54}, Cambridge University Press, 1995.

\end{thebibliography}

\end{document}